\subsection{Criba de Conteo de Divisores}

\graybold{Preguntas guía:}

\begin{itemize}
\item ¿Necesito el número de divisores de MUCHOS números distintos, todos acotados por un mismo límite razonable ($\le 10^6$-$10^7$)?
\item ¿Factorizar cada número por separado (aunque sea con Pollard's Rho) sería un desperdicio si hay muchas consultas repetidas o cercanas en rango?
\item ¿Puedo precalcular la respuesta para TODOS los valores posibles de una sola vez?
\end{itemize}

\graybold{¿Para qué sirve?} Precalcular, para todos los enteros hasta un límite A, su cantidad de divisores en una sola criba, en vez de factorizar cada consulta por separado. Rentable cuando hay muchas consultas (n grande) sobre un rango de valores acotado.

\graybold{Complejidad:} \bigO{A log A} para construir la criba; \bigO{1} por consulta.

\graybold{Fundamento teórico:} En vez de factorizar cada número, se recorre cada posible divisor $i$ de 1 a A, y se le suma 1 a la cuenta de TODOS sus múltiplos (ya que $i$ es divisor de $i, 2i, 3i, \ldots$). Al terminar, \texttt{divcount[x]} contiene la cantidad exacta de divisores de x. Esto es análogo a la Criba de Eratóstenes, pero contando divisores en vez de marcar primos; se vectoriza con numpy asignando por slicing (\texttt{divcount[i::i] += 1}) en vez de un loop interno en Python puro.

\graybold{Ejercicio aplicado}

Dados n enteros (cada uno hasta 10\textsuperscript{6}), reportar la cantidad de divisores de cada uno.

\graybold{I/O:} n ≤ 10\textsuperscript{5}, valores ≤ 10\textsuperscript{6} → lectura total con tokenización (patrón 2); la criba se vectoriza con numpy (\textasciitilde{}1.7s para construirla una sola vez, independientemente de n).

\graybold{Código solución}

\begin{lstlisting}[language=Python]
import sys
import numpy as np

def solve():
    data = sys.stdin.buffer.read().split()
    n = int(data[0])
    xs = list(map(int, data[1:1 + n]))
    maxv = max(xs) if xs else 1

    divcount = np.zeros(maxv + 1, dtype=np.int32)
    for i in range(1, maxv + 1):
        divcount[i::i] += 1        # cada multiplo de i tiene a i como divisor

    out = [str(divcount[x]) for x in xs]
    sys.stdout.write("\n".join(out) + "\n")

solve()
\end{lstlisting}